\section{Monte Carlo：从随机样本到数值答案}

许多积分和期望无法解析计算，却可以从相应分布中抽样。Monte Carlo 方法把数值问题改写成期望，再用样本平均估计：

\[
I=\E[g(X)]
\approx
\widehat I_n
=\frac1n\sum_{i=1}^{n}g(X_i).
\]
这里 $X_i$ 是从目标分布独立抽取的样本。

\subsection{大数定律提供一致性}

若

\[
\E[|g(X)|]<\infty,
\]
则

\[
\widehat I_n\to I
\]
依概率，适当条件下几乎处处。因此样本越多，估计最终会稳定到目标期望。

这解释了方法为什么成立，却没有给出有限样本误差。

\subsection{中心极限定理给出标准误}

若

\[
\sigma_g^2=\operatorname{Var}(g(X))<\infty,
\]
则

\[
\sqrt n(\widehat I_n-I)
\xrightarrow{d}
\mathcal N(0,\sigma_g^2).
\]
所以估计的标准误约为

\[
\frac{\sigma_g}{\sqrt n}.
\]
用样本方差

\[
s_g^2
=\frac1{n-1}\sum_{i=1}^{n}
\bigl(g(X_i)-\widehat I_n\bigr)^2
\]
估计 $\sigma_g^2$，便得到

\[
\operatorname{SE}(\widehat I_n)
\approx\frac{s_g}{\sqrt n}.
\]
一个常用近似 (95%) 误差区间是

\[
\widehat I_n
\pm1.96\frac{s_g}{\sqrt n}.
\]
要把标准误缩小一半，样本量要增加到四倍。这种 $1/\sqrt n$ 收敛较慢，却在维度上不像规则网格积分那样直接遭受 $m^d$ 个网格点的爆炸。

\subsection{用随机点估计 $\pi$}

在正方形 $[-1,1]^2$ 中均匀抽点 ($X_i,Y_i$)，定义

\[
Z_i=\mathbf1_{\{X_i^2+Y_i^2\le1\}}.
\]
圆面积占正方形面积比例为

\[
P(Z_i=1)=\frac{\pi}{4}.
\]
因此

\[
\widehat\pi_n
=4\bar Z_n
\]
是 $\pi$ 的无偏估计。

因为 $Z_i$ 是 Bernoulli，

\[
\operatorname{Var}(\widehat\pi_n)
=\frac{16p(1-p)}n,
\qquad p=\frac\pi4.
\]
这个例子直观，却不是计算 $\pi$ 的高效方法。它的教学价值在于展示怎样把几何面积变成事件概率；它也反向使用了``大数定律与长期平均''一节建立的同一联系：那里由模型概率 $p$ 预言长期频率，这里由长期频率反推未知的 $p$。

\subsection{方差缩减比盲目加样本更聪明}

Monte Carlo 误差常数由 $\operatorname{Var}(g(X))$ 决定。除了增大 $n$，还可以改变估计器而保持目标期望：

\begin{itemize}
\tightlist
\item
  \textbf{控制变量}：利用一个期望已知且与目标相关的变量抵消波动；
\item
  \textbf{对偶变量}：让成对样本产生负相关误差；
\item
  \textbf{分层抽样}：保证重要子区域都得到样本；
\item
  \textbf{重要性抽样}：从更常访问重要区域的分布抽样，再用权重修正。
\end{itemize}

重要性抽样把

\[
I=\int g(x)p(x)\,dx
\]
写成

\[
I=\E_q\left[
g(X)\frac{p(X)}{q(X)}
\right],
\]
其中下标 $q$ 表示 $X$ 改从密度 $q$ 抽取，前提是 $q(x)>0$ 覆盖所有 $g(x)p(x)\ne0$ 的区域。若 $q$ 选得差，权重可能重尾，方差甚至无穷。

\subsection{独立样本是假设，不是界面按钮}

伪随机数生成器要有足够质量；马尔可夫链 Monte Carlo（MCMC）样本通常相关，不能把样本数 $n$ 当作 $n$ 个独立观察。相关性会降低有效样本量：

\[
n_{\mathrm{eff}}<n.
\]
此外还要区分：

\begin{itemize}
\tightlist
\item
  \textbf{Monte Carlo 随机误差}：可通过重复与标准误评估；
\item
  \textbf{离散化偏差}：例如数值求解近似模型带来的系统误差；
\item
  \textbf{模型偏差}：目标分布本身不代表现实；
\item
  \textbf{初始与收敛偏差}：MCMC 尚未进入平稳状态。
\end{itemize}

样本平均非常稳定地收敛到错误目标，仍然是错误答案——这与``大数定律与长期平均''一节列出的选择偏差、共同系统误差同源：重复只能缩小随机误差，不能修正抽样机制本身。

\subsection{一份可报告的 Monte Carlo 结果}

至少应说明：

\begin{enumerate}
\def\labelenumi{\arabic{enumi}.}
\tightlist
\item
  目标期望怎样定义；
\item
  样本从什么分布、用什么机制产生；
\item
  样本量与是否独立；
\item
  点估计；
\item
  标准误或可解释的误差界；
\item
  是否使用方差缩减；
\item
  可能存在的系统偏差。
\end{enumerate}

随机计算不是``多跑几次取平均''这么简单。它是一套把概率极限定理、抽样设计和误差报告接成闭环的方法。

\subsubsection{自己补完}

要估计

\[
I=\int_0^1 e^{-x^2}\,dx,
\]
令 $U_i\sim\mathrm{Unif}(0,1)$，并取

\[
\widehat I_n=\frac1n\sum_{i=1}^{n}e^{-U_i^2}.
\]
\begin{enumerate}
\def\labelenumi{\arabic{enumi}.}
\tightlist
\item
  证明估计目标确实是 $I$；
\item
  写出估计标准误的方法；
\item
  若希望标准误缩小到原来的三分之一，样本量需放大多少倍？
\item
  设计一个简单分层方案，把 $[0,1]$ 分成两段分别抽样。
\end{enumerate}
